\begin{casebox}
\textbf{知识点小案例：atomic flag 不会自动发布普通字段。}
特例是生产者先写 \code{config.worker\_count}，再执行 \code{ready.store(true, relaxed)}；消费者看到 \code{ready.load(relaxed)} 为 true 后读取 config。这个程序的 atomic flag 自身没有数据竞争，但它没有建立“配置字段先于读取方可见”的发布关系。改法是生产者用 release store，消费者用 acquire load，并且 config 在发布后不再被无保护修改。由此推出规律：内存序要从用途推导，指标计数可以 relaxed，发布普通数据需要 release/acquire。工程上每个 atomic 都要写变量用途、发布数据、读取方动作和生命周期规则。
\end{casebox}
